\documentclass[a4paper]{article}

\usepackage{graphicx}
\usepackage{amsmath,amssymb}
\usepackage{minted}
\usepackage{multirow}
\usepackage{float}
\usepackage{todonotes}
\usepackage{forest}
\usepackage{xstring}
\usepackage{xcolor}
\usepackage[most]{tcolorbox}
\usepackage{xparse}
\usepackage{natbib}

\usetikzlibrary{graphs,graphdrawing}
\usegdlibrary{layered}

% Benötigt: \usepackage{xcolor}
\definecolor{TranslationGray}{RGB}{120,120,120}
\newcommand{\EN}[1]{\par{\small\textcolor{TranslationGray}{#1}}\vspace{0.4em}}

% for external graphviz
\input{/home/donkarlo/Dropbox/repo/nd_language_ai_project/src/nd_language_ai/formal/computer/programming/tex/latex/code_clips/external_graphviz.tex}

% for tags
\input{/home/donkarlo/Dropbox/repo/nd_language_ai_project/src/nd_language_ai/formal/computer/programming/tex/latex/parts/control_sequence/kind/semtag/semtag.tex}

% for links
\input{/home/donkarlo/Dropbox/repo/nd_language_ai_project/src/nd_language_ai/formal/computer/programming/tex/latex/code_clips/seehere.tex}



\title{Kasus}
\date{}

\begin{document}
    \maketitle
    
    
    \section{Kasus bei deutschen Verben}
        
        \textbf{Im Deutschen reicht es nicht immer, nur direktes und indirektes Objekt zu erkennen.}
        \EN{In German, it is not always enough to recognize only direct and indirect objects.}
        
        \textbf{Oft bestimmt das Verb selbst, welchen Kasus das Objekt bekommt.}
        \EN{Often the verb itself determines which case the object takes.}
        
        \textbf{Deshalb sollte man viele Verben zusammen mit ihrem Kasus lernen.}
        \EN{Therefore, many verbs should be learned together with their case.}
        
        \subsection{Verben mit Dativ}
            
            \textbf{Einige Verben brauchen immer oder meistens den Dativ.}
            \EN{Some verbs always or usually require the dative.}
            
            \textbf{helfen + Dativ}
            \EN{to help + dative}
            
            \textbf{Ich helfe dir.}
            \EN{I help you.}
            
            \textbf{danken + Dativ}
            \EN{to thank + dative}
            
            \textbf{Ich danke dem Großvater.}
            \EN{I thank the grandfather.}
            
            \textbf{antworten + Dativ}
            \EN{to answer + dative}
            
            \textbf{Ich antworte der Lehrerin.}
            \EN{I answer the teacher.}
            
            \textbf{gehören + Dativ}
            \EN{to belong to + dative}
            
            \textbf{Das Buch gehört mir.}
            \EN{The book belongs to me.}
        
        \subsection{Verben mit Akkusativ}
            
            \textbf{Die meisten Verben mit Objekt brauchen den Akkusativ.}
            \EN{Most verbs with an object require the accusative.}
            
            \textbf{sehen + Akkusativ}
            \EN{to see + accusative}
            
            \textbf{Ich sehe dich.}
            \EN{I see you.}
            
            \textbf{haben + Akkusativ}
            \EN{to have + accusative}
            
            \textbf{Ich habe einen Kaffee.}
            \EN{I have a coffee.}
            
            \textbf{kaufen + Akkusativ}
            \EN{to buy + accusative}
            
            \textbf{Ich kaufe das Brot.}
            \EN{I buy the bread.}
        
        \subsection{Verben mit zwei Objekten}
            
            \textbf{Wenn ein Satz zwei Objekte hat, ist die Person oft Dativ und die Sache Akkusativ.}
            \EN{When a sentence has two objects, the person is often dative and the thing is accusative.}
            
            \textbf{Ich gebe dir das Buch.}
            \EN{I give you the book.}
            
            \textbf{dir = Dativ}
            \EN{to you = dative}
            
            \textbf{das Buch = Akkusativ}
            \EN{the book = accusative}
            
            \textbf{Ich werfe dir den Ball zu.}
            \EN{I throw the ball to you.}
            
            \textbf{dir = Dativ}
            \EN{to you / receiver = dative}
            
            \textbf{den Ball = Akkusativ}
            \EN{the ball = accusative}
        
        \subsection{Wichtige Regel}
            
            \textbf{Für Deutsch muss man zuerst lernen, welchen Kasus ein Verb verlangt.}
            \EN{For German, one must first learn which case a verb requires.}
            
            \textbf{Direktes und indirektes Objekt helfen beim Verstehen, aber sie reichen nicht immer aus.}
            \EN{Direct and indirect objects help with understanding, but they are not always enough.}
            
            \textbf{Die beste Methode ist: Verb + Kasus + Beispiel lernen.}
            \EN{The best method is: learn verb + case + example.}
            
            \textbf{helfen + Dativ: Ich helfe dir.}
            \EN{to help + dative: I help you.}
            
            \textbf{sehen + Akkusativ: Ich sehe dich.}
            \EN{to see + accusative: I see you.}
            
            \textbf{geben + Dativ + Akkusativ: Ich gebe dir das Buch.}
            \EN{to give + dative + accusative: I give you the book.}
            
            \nocite{*}
%            \bibliographystyle{plainnat}
%            \bibliography{/home/donkarlo/Dropbox/repo/nd_shared_research/bibliography/bibliography.bib}
\end{document}